\section{Table-header, continuous and shared legends}
\subsection{Horizontal categories and combined headings}
\key{\gradtablelegend} takes optional graphic settings and a list of names,
drawing a horizontal category row at its natural width.
\key{\gradtablelegendhead} also takes a title, centered above the row with a
2pt gap. Both commands belong to core and
work with \key{libraries=none}; the optional \key{legend} library is unnecessary.
They render only when explicitly called, without inserting legends into other
graphics or tables.
\begin{demo}{Service response: a heading with five status items}
\begingroup
\gradbarssetup{stack colors={gray,teal,blue,
  orange,red}}
\begin{center}
\begin{tabular}{lc}
Service & \gradtablelegendhead{Response status}{Pending,Fast,Normal,Slow,Timeout}\\
\hline
Lookup & \gradbar[max=1000,width=42mm,color=teal,precision=0,unit={ ms}]{120}\\
Export & \gradbar[max=1000,width=42mm,color=blue,precision=0,unit={ ms}]{420}\\
\end{tabular}
\par\medskip
\gradtablelegendhead[table legend=false]{Response status}
  {Pending,Fast,Normal,Slow,Timeout}
\end{center}
\endgroup
\end{demo}
These are \key{/gradbars} graphic keys for helper options or
\key{\gradbarssetup}, with normal grouping and left-to-right overrides;
they are not new native \key{gradtable} declaration keys.
\key{table legend=true|false} defaults to true; the bare key means true.
False hides the category row; the heading helper retains its title without
reserving the 2pt legend gap. \key{legend font} defaults to
\key{\sffamily\scriptsize} for category labels.
\key{legend swatch size=6pt} must be positive;
\key{legend text gap=2pt} and \key{legend item gap=4pt} must be nonnegative.
The border defaults are \key{legend border color=black} and
\key{legend border width=.4pt}; border width must be nonnegative.

Registered \key{\gradbarscategory} mappings take precedence; other names use
positional \key{stack colors}, cycling as needed. Print mode reuses category
textures. The helpers do not infer thresholds, classify observations or change
bar colors, so the example sets bar colors explicitly. The older vertical
\key{\gradbarslegend} and stack \key{legend=true} remain independent of
\key{table legend}.
Use the heading helper directly in a normal table header or a \key{\gradcolumn}
title. In a collected numeric column, bypass numeric collection with:
\begin{source}
\gradbarshead{\gradtablelegendhead{Response status}
  {Pending,Fast,Normal,Slow,Timeout}}
\end{source}
It replaces handwritten title/minipage/TikZ-swatch macros without inserting
\key{\multicolumn}; retain any column spanning in the surrounding table code.
See \key{docs/examples/service-response-legend.tex} for a standalone example.
The category graphic follows existing \key{alt}, \key{accessibility} and
\key{category-legend} backend rules, including explicit native fallback when
required. The title remains ordinary text outside the legend Figure.

\subsection{Continuous and shared legends}
Load \key{legend} in the preamble; it depends on Heat. Unlike category swatches,
a continuous legend describes a numeric scale and uses the same color interpolation.
\begin{demo}{Two values sharing a frozen scale and legend}
\begingroup
\gradscale{legend-score}{min=-10,max=30}
\gradlegenddefine{legend-score}{scale group=legend-score,
  heat palette=purple-teal,
  heat center=0,legend title={Score}}
\gradheat[legend group=legend-score]{-5}
\gradheat[legend group=legend-score]{20}
\par\medskip
\gradlegenduse[width=60mm,height=3mm,legend missing=false]{legend-score}
\endgroup
\end{demo}
\key{\gradheatlegend} accepts ordinary graphic keys plus \key{legend title},
\key{legend ticks}.\par
\key{legend missing=false} hides the missing-value item and its automatic description;
\key{true} is the default. Both legend commands accept this switch.\par
Set \key{legend orientation} to \key{horizontal} or \key{vertical}.
Width is ramp length; height is thickness. Endpoints and an interior neutral
center are default ticks; custom ticks must increase within the range. The ramp
uses 256 vector strips sampled from Heat, with decorative semitransparent white
separators in both orientations and print mode. These lines change the texture,
not the numeric scale; they are not category or bin boundaries.
Increase precision or spacing if tick text overlaps.
\key{\gradlegenddefine} freezes scale, direction, unit and color roles in the
current group. Use \key{legend group=name} in standalone graphics or column
options. Unknown names, duplicate declarations and contradictory scale/color/unit
settings are errors. \key{\gradlegenduse} only accepts layout, tick formatting,
font and accessibility overrides, leaving subsequent graphics unchanged.
\begin{demo}{Reference symbols have different statistical meanings}
\gradreferencelegend{target/Declared target,
  reference/Reference value,band/Reference range,
  error/Error distances,box/Quartiles and observed whiskers}
\end{demo}
Reference descriptions declare their actual meaning; neither error distances
nor box whiskers automatically denote confidence intervals.
For shared category or reference legends, set \key{legend kind=categorical|reference}
and \key{legend names}. Repeated decoration may use \key{accessibility=artifact}.


\subsection{Built-in Heat palettes}
One \key{heat palette=name} selects endpoint colors, a white center and the interpolation mode.
Sequential palettes \key{blue}, \key{teal}, \key{purple} and \key{orange} run from white to that color.
Diverging palettes \key{purple-teal}, \key{blue-red}, \key{orange-blue} and \key{brown-teal}
run from the first color through white to the second; \key{heat center} still sets the numeric center.
These are color shortcuts, not accessibility certifications; existing accessible palettes remain available.
Explicit low, center and high colors placed after the preset still override it.
\begin{demo}{Built-in continuous palettes}
\gradheatlegend[width=45mm,height=2mm,min=-10,max=30,heat center=0,
  heat palette=blue,legend title={blue},legend missing=false]\quad
\gradheatlegend[width=45mm,height=2mm,min=-10,max=30,heat center=0,
  heat palette=teal,legend title={teal},legend missing=false]\par\medskip
\gradheatlegend[width=45mm,height=2mm,min=-10,max=30,heat center=0,
  heat palette=purple,legend title={purple},legend missing=false]\quad
\gradheatlegend[width=45mm,height=2mm,min=-10,max=30,heat center=0,
  heat palette=orange,legend title={orange},legend missing=false]\par\medskip
\gradheatlegend[width=45mm,height=2mm,min=-10,max=30,heat center=0,
  heat palette=purple-teal,legend title={purple-teal},legend missing=false]\quad
\gradheatlegend[width=45mm,height=2mm,min=-10,max=30,heat center=0,
  heat palette=blue-red,legend title={blue-red},legend missing=false]\par\medskip
\gradheatlegend[width=45mm,height=2mm,min=-10,max=30,heat center=0,
  heat palette=orange-blue,legend title={orange-blue},legend missing=false]\quad
\gradheatlegend[width=45mm,height=2mm,min=-10,max=30,heat center=0,
  heat palette=brown-teal,legend title={brown-teal},legend missing=false]
\end{demo}
